\chapter{Mengendalikan Sistem}
\label{ch:kendali}

Di semester pertama saya mengambil dua mata kuliah yang sekilas tidak ada hubungannya dengan
machine learning. Control systems membahas sistem dinamis linear dari model ruang keadaan sampai
pengendali diskret, dengan latihan dan ujian berbasis MATLAB. Production automation systems membahas
dunia pabrik: perangkat otomasi, sistem waktu nyata, pemrograman PLC, fieldbus, dan standar IEC
61499 untuk sistem otomasi terdistribusi. Belakangan terlihat bahwa teori kendali adalah saudara
kandung reinforcement learning, dan bahwa model ruang keadaan adalah bahasa yang sama yang dipakai
Kalman filter, state space model seperti Mamba, dan penemuan persamaan dari data.

\section{Sistem dan model}

Sebuah sistem dinamis dijelaskan oleh keadaannya, yaitu sekumpulan besaran yang cukup untuk
menentukan masa depannya kalau input diketahui. Untuk rangkaian listrik RC, keadaannya tegangan
kapasitor. Untuk kaskade tangki air yang saling mengalirkan, keadaannya ketinggian air di setiap
tangki. Untuk populasi pemangsa dan mangsa, keadaannya jumlah masing-masing hewan, dan persamaan
Lotka--Volterra menggambarkan siklus naik turun yang khas: mangsa berlimpah membuat pemangsa
bertambah, pemangsa yang banyak membuat mangsa berkurang, dan seterusnya. Bentuk umumnya adalah
\begin{equation}
\dot{\vx}=\bm f(\vx,\bm u),\qquad \vy=\bm h(\vx,\bm u),
\end{equation}
dengan $\bm u$ input yang bisa kita atur dan $\vy$ keluaran yang bisa kita ukur.

Sebagian besar teori dibangun untuk sistem linear yang tidak berubah terhadap waktu,
\begin{equation}
\dot{\vx}=\mA\vx+\bm B\bm u,\qquad \vy=\bm C\vx+\bm D\bm u.
\end{equation}
Sistem nonlinear dilinearkan di sekitar titik kesetimbangan dengan mengganti $\bm f$ dengan
matriks Jacobiannya. Solusi sistem linear adalah
$\vx(t)=e^{\mA t}\vx_0+\int_0^te^{\mA(t-\tau)}\bm B\bm u(\tau)\dd\tau$, dan perilakunya ditentukan oleh nilai
eigen $\mA$. Sistem stabil secara asimtotik, artinya kembali ke kesetimbangan dari gangguan apa pun,
tepat ketika semua nilai eigen punya bagian real negatif. Bentuk Jordan menguraikan sistem menjadi
mode-mode independen, sehingga setiap nilai eigen bisa dibaca sebagai satu cara sistem meluruh,
tumbuh, atau berosilasi.

\section{Bisa dicapai, bisa diamati}

Dua pertanyaan menentukan apakah sebuah sistem bisa dikendalikan dengan baik. Pertama, apakah input
bisa membawa keadaan ke mana saja? Untuk sistem linear, jawabannya ya tepat ketika matriks
\begin{equation}
\mathcal R=\big[\bm B\ \ \mA\bm B\ \ \mA^2\bm B\ \cdots\ \mA^{n-1}\bm B\big]
\end{equation}
berpangkat penuh. Kedua, apakah keadaan bisa direkonstruksi dari keluaran? Jawabannya ya tepat ketika
matriks yang disusun dari $\bm C,\bm C\mA,\dots,\bm C\mA^{n-1}$ berpangkat penuh. Dekomposisi Kalman
memisahkan keadaan menjadi bagian yang bisa dicapai dan diamati, bagian yang hanya salah satunya,
dan bagian yang tidak keduanya. Bagian yang tidak terjangkau dan tidak stabil adalah kabar buruk,
karena tidak ada pengendali yang bisa menyelamatkannya.

\section{Pengendali dan pengamat}

Dengan umpan balik keadaan $\bm u=-\bm K\vx$, sistem tertutup menjadi $\dot{\vx}=(\mA-\bm B\bm K)\vx$, dan
kita bisa memilih $\bm K$ agar nilai eigen sistem tertutup berada di tempat yang kita inginkan, selama
sistemnya bisa dicapai. Ambil contoh paling sederhana, sebuah massa tanpa gesekan, $\ddot x=u$. Dengan
keadaan posisi dan kecepatan serta $u=-k_1x-k_2\dot x$, persamaan karakteristiknya $s^2+k_2s+k_1=0$.
Kalau kita menginginkan kedua kutub di $-2$, polinomialnya harus $(s+2)^2=s^2+4s+4$, sehingga
$k_1=k_2=4$. Rumus Ackermann melakukan perhitungan yang sama untuk sistem dengan satu input
berukuran berapa pun.

Kalau keadaan tidak diukur seluruhnya, kita membangun pengamat yang menirukan sistem dan mengoreksi
dirinya dengan galat keluaran,
\begin{equation}
\dot{\hat\vx}=\mA\hat\vx+\bm B\bm u+\bm L\big(\vy-\bm C\hat\vx\big).
\end{equation}
Prinsip pemisahan menyatakan bahwa pengendali dan pengamat bisa dirancang terpisah lalu digabung,
dan kutub-kutub sistem gabungan adalah gabungan kutub keduanya.

Alih-alih memilih kutub, kita juga bisa meminta pengendali terbaik menurut sebuah biaya. Linear
quadratic regulator meminimalkan
\begin{equation}
J=\int_0^\infty\big(\vx^\top\bm Q\vx+\bm u^\top\bm R\bm u\big)\dd t,
\end{equation}
yang menimbang seberapa jauh keadaan menyimpang terhadap seberapa besar tenaga kendali yang
dipakai. Jawabannya tetap berupa umpan balik linear, dengan $\bm K$ dari persamaan Riccati. Ketika
pengukurannya berisik, pengamat optimalnya adalah Kalman filter \citep{kalman1960}, dan gabungan
LQR dengan Kalman filter disebut LQG. Dari sudut pandang \cref{ch:rl}, LQR adalah masalah
reinforcement learning yang modelnya diketahui sempurna, dinamikanya linear, dan biayanya kuadratik,
sehingga persamaan Bellman bisa diselesaikan secara tertutup.

\section{Domain frekuensi}

Transformasi Laplace mengubah persamaan diferensial linear menjadi persamaan aljabar. Hubungan
antara input dan keluaran menjadi fungsi alih, $G(s)=\bm C(s\bm I-\mA)^{-1}\bm B+\bm D$. Sistem stabil
dalam arti BIBO, yaitu input terbatas selalu menghasilkan keluaran terbatas, kalau semua kutub $G(s)$
berada di setengah bidang kiri. Respons frekuensi $G(i\omega)$ menyatakan bagaimana sinyal sinus
diperkuat dan digeser fasenya, dan diagram Bode menggambarnya dalam skala logaritmik.

Dalam lingkar umpan balik, pengendali $R(s)$ dan proses $G(s)$ membentuk fungsi lingkar terbuka
$L(s)=R(s)G(s)$, dan banyak sifat sistem tertutup bisa dibaca dari bentuk $L$. Pengendali
proporsional-integral menambahkan suku integral yang menghapus galat keadaan tunak, karena selama
galat masih ada, integralnya terus tumbuh dan mendorong input. Kuliah merancang pengendali seperti ini
dengan membentuk respons langkah yang diinginkan, lalu memeriksa kestabilannya dengan kriteria
grafis. Bagian terakhir kuliah membawa semuanya ke dunia digital. Pengendali di komputer bekerja
pada sampel, dan pemetaan $z=e^{sT}$ menghubungkan domain kontinu dengan domain diskret. Transformasi
Tustin memberi cara praktis untuk mengubah pengendali kontinu menjadi algoritme diskret.

\section{Otomasi di pabrik}

Teori kendali bertemu kenyataan di lantai pabrik. Kuliah production automation memperlihatkan
lapisan-lapisan yang menjalankan kendali itu. \istilah{Programmable logic controller} (PLC) membaca
semua input, menjalankan programnya, lalu menulis semua output, dalam siklus yang berulang setiap
beberapa milidetik. Programnya ditulis dalam bahasa standar IEC 61131-3, seperti diagram tangga yang
meniru rangkaian relai, diagram blok fungsi, atau teks terstruktur. Sistem waktu nyata membedakan
tenggat keras, yang tidak boleh terlewat sama sekali karena akibatnya berbahaya, dari tenggat lunak,
yang sesekali boleh terlewat. Fieldbus menghubungkan sensor, aktuator, dan pengendali dengan
komunikasi yang waktunya bisa dipastikan. Standar IEC 61499 memecah aplikasi otomasi menjadi blok fungsi
yang digerakkan oleh kejadian dan bisa disebar ke banyak perangkat, sebuah langkah menuju pabrik yang
bisa dikonfigurasi ulang dengan cepat, tema utama gerakan Industri 4.0.

\section{Kendali dan machine learning}

Hubungan antara kendali dan machine learning berjalan ke dua arah. Dari kendali ke machine learning:
model ruang keadaan menjadi state space model di deep learning (\cref{ch:lstm}), metode adjoint dari
kendali optimal menjadi cara melatih neural ODE (\cref{ch:mlat}), dan Kalman filter menjadi kerangka
data assimilation (\cref{ch:prob}). Dari machine learning ke kendali: model dinamika bisa dipelajari
dari data alih-alih diturunkan dari hukum fisika, lalu dipakai dalam pengendali prediktif yang
mengoptimalkan beberapa langkah ke depan di setiap saat. Teori operator Koopman \citep{koopman1931}
menawarkan cara melihat sistem nonlinear sebagai sistem linear di ruang fungsi yang lebih besar,
sehingga alat-alat kendali linear bisa dipakai kembali, sebuah gagasan yang kita temui lagi di
\cref{ch:penemuan}. Reinforcement learning sendiri bisa dipandang sebagai kendali optimal tanpa model,
dan pada sistem nyata keduanya sering digabung: model fisika memberi dasar, pembelajaran mengoreksi
apa yang tidak dimodelkan.

\sumberbab{Bab ini mengikuti naskah kuliah Control Systems (sistem dan model, sifat sistem, sistem LTI,
keterjangkauan dan keteramatan, pengendali dan pengamat keadaan, LQR dan LQG, transformasi Laplace,
desain pengendali, sistem diskret) dan kuliah Production Automation Systems (perangkat otomasi, sistem
waktu nyata, PLC, fieldbus, IEC 61499). Bacaan lanjut tentang hubungan kendali dan data:
\citet{brunton2022}.}
